% Appendices: source-aligned review, 3 October 2026.
\appendix
\chapter{LFS Questionnaire Fields}
\label{app:questionnaire}

Table~\ref{tab:app-questionnaire} documents selected fields and canonical
values from \texttt{ConversationManager}'s correction schema and dynamic
field order. The item codes are implementation labels, not evidence that
the prototype reproduces a complete official UAE questionnaire. Interview
prompts vary by language and state. The schema describes intended values,
not proof that every translated answer is normalized correctly.

\begingroup
\setlength{\tabcolsep}{3pt}
\footnotesize
\renewcommand{\arraystretch}{1.35}
\begin{longtable}{@{}>{\raggedright\arraybackslash}p{0.06\textwidth}>{\raggedright\arraybackslash}p{0.27\textwidth}
                       >{\raggedright\arraybackslash}p{0.36\textwidth}>{\raggedright\arraybackslash}p{0.23\textwidth}@{}}
\caption{Selected implemented fields, canonical values, and routing gates.}
\label{tab:app-questionnaire}\\
\toprule
\textbf{Code} & \textbf{Field} & \textbf{Canonical values or input} &
\textbf{Routing} \\
\midrule
\endfirsthead
\multicolumn{4}{c}{\tablename\ \thetable{} --- continued} \\
\toprule
\textbf{Code} & \textbf{Field} & \textbf{Canonical values or input} &
\textbf{Routing} \\
\midrule
\endhead
\midrule
\multicolumn{4}{r}{\small Continued on the next page} \\
\endfoot
\bottomrule
\endlastfoot
C1 & \texttt{employment\_status} &
  \texttt{employed}, \texttt{unemployed},
  \texttt{not\_in\_labour\_force} & All paths \\
C3 & \texttt{employment\_nature} &
  \texttt{paid\_employee}, \texttt{employer}, \texttt{self\_employed},
  \texttt{contributing\_family\_member} & Employed \\
C5 & \texttt{job\_title} & Free-text current occupation &
  Employed; ISCO-08 input \\
C5a & \texttt{job\_duties} & Free-text duties &
  Employed; occupation-refinement input \\
C6 & \texttt{industry} & Free-text economic activity &
  Employed; ISIC Rev.~4 input \\
D7 & \texttt{contract\_type} &
  \texttt{permanent}, \texttt{temporary}, \texttt{no\_contract} &
  Employed paid employees \\
E1 & \texttt{monthly\_wage\_range} &
  AED bands: \texttt{under\_5000}, \texttt{5000\_to\_10000},
  \texttt{10000\_to\_20000}, \texttt{20000\_to\_30000},
  \texttt{above\_30000} & Employed paid employees \\
F1 & \texttt{job\_search\_active} & \texttt{yes}, \texttt{no} &
  Unemployed \\
F3 & \texttt{available\_for\_work} &
  \texttt{yes}, \texttt{no}; availability within two weeks &
  Unemployed; F1=no and F3=no changes status to outside the labour force \\
B6 & \texttt{field\_of\_study} & Free-text main field &
  All paths with education \texttt{bachelor}, \texttt{master}, or \texttt{phd} \\
\end{longtable}
\endgroup

Dynamic order also depends on nationality, secondary employment, platform
work, and previous employment (Section~\ref{sec:arch-fsm}). The wage gate
belongs to the conversation manager, not the ValidationAgent's rules.
A known extraction defect can populate wages from an unrelated numeric
answer and skip the later question. Hindi and Urdu status options also
have a normalization defect. The table records intended routing rather
than an assertion that every path is correct.

\chapter{Agent Configuration Reference}
\label{app:crewai}

Agents are constructed in Python. These listings are descriptive
configuration summaries, not executable YAML files. Interactive defaults
are separated from optional evaluation features.

\section*{B.1 --- Occupation Classifier}
\begin{lstlisting}[caption={Current occupation-classifier configuration summary.}]
interactive_route: ISCOClassifier()
catalogue_profile: legacy default
retrieval: hierarchical default; flat available by explicit option
encoder: multilingual-e5-small default
official_profiles:
  unit_groups: 436
  encoders: multilingual-e5-small or multilingual-e5-large
  text: title-only or enriched definitions/tasks/examples
hitl_threshold: 0.70
skip_rerank_threshold: 0.92
reranking:
  default_factory: TaskType.GENERAL
  explicit_model: get_llm_strict(model, temperature)
  enable_llm_false: retrieval only
  unavailable_default_agent: retrieval-only degradation
\end{lstlisting}

The official enriched large-encoder flat profile supplies the principal
full-heldout retrieval result, not the live interview default. GENERAL routing
prefers a configured local Ollama model and can fall back through cloud
providers. CRITICAL routing selects the configured Claude model without
the GENERAL fallback chain. The current occupation constructor itself
selects GENERAL when unpinned, despite its older module documentation.

\section*{B.2 --- Semantic Relation Engine}
\begin{lstlisting}[caption={Implemented coherence rules and score settings.}]
source: project-authored plausibility tables
inputs: ISCO-08 code, ISIC section, ISCED 2011 level
not_an_input: ISCED-F 2013 detailed field
raw_score:
  both_dimensions: 0.55 * industry_ok + 0.45 * education_ok
  one_dimension: available dimension's compatibility score
  neither_dimension: 0.80
partial_credit: 0.15 * moderate_count * (1 - raw_score)
score: clipped to [0, 1], rounded to 4 decimal places
coherent_threshold: 0.70
education_gap_severity:
  1: LOW
  2: MODERATE
  3_or_more: HIGH
confidence_adjustment:
  score_ge_0.90: +0.10
  score_ge_0.70: +0.05
  score_ge_0.50: -0.05
  otherwise: -0.20
education_submajor_minimums: {"22": 7, "23": 6, "25": 6, "26": 7}
industry_submajor_examples: {"22": ["Q"], "23": ["P"], "25": ["J"]}
\end{lstlisting}

Education severity uses the distance outside the allowed level range.
Industry severity uses a different rule: an industry allowed at the major
level but excluded by a stricter sub-major rule is MODERATE; an industry
outside the applicable allowed lists is HIGH. Sub-major 11 has no special
education override and inherits the managers' major-group range.
Education rules concern expected attainment, not legally required
qualifications or field-of-study correspondence.

Returning an adjustment value is distinct from applying it to a saved
classification. These hand-set rules and thresholds do not produce an
empirically calibrated probability or an official ILO/UNESCO crosswalk.

\section*{B.3 --- Validation Agent}
\begin{lstlisting}[caption={Validation rules and semantic stage.}]
R01: employment_status and industry required for all statuses;
     job_title, hours_per_week, employment_type also required if employed
R02: supplied hours_per_week must be numeric
R03: hours_per_week must be within [1, 168]
R04: hours_per_week > 80 gives an outlier warning
R05: unemployed with current job_title gives a contradiction
R06: unemployed with hours_per_week > 0 gives a contradiction
R07: outside labour force with current job_title gives a contradiction
R08: full_time with hours_per_week < 20 gives a warning
R09: part_time with hours_per_week >= 35 gives a warning
R10: employed with hours_per_week = 0 gives an error
semantic_stage: attempted if deterministic checks contain no errors
semantic_model_factory: TaskType.CRITICAL
\end{lstlisting}

The route calls \texttt{validate()} without explicitly disabling its
semantic stage, and catches validation failures. A successful result is
not required for completion. Required fields also need reconciliation
with inactive paths (Section~\ref{sec:impl-validation}).

\chapter{Semantic Relation Engine: Validation Evidence}
\label{app:sre-demo}

The available evidence checks implementation consistency and escalation.
The 61-case fixture in \path{eval/sre_expanded_validation.py} assigns
expected outcomes using the same project's plausibility rules. Agreement
on all 61 cases supports execution of those rules; it does not
independently validate their accuracy on labour-force data.

The endpoint script reuses this fixture and checks for queue entries
tagged as HIGH semantic incoherence. Three documented executions each
recorded 15 of 15 expected HIGH cases queued and zero of 46 other cases
queued for that reason:
\begin{table}[htbp]
\centering
\caption{Repeated endpoint checks on the same 61-case fixture.}
\label{tab:app-sre-runs}
\begin{tabular}{@{}l r r@{}}
\toprule
\textbf{Execution} & \textbf{Expected HIGH queued} &
\textbf{Other cases with SRE HIGH entry} \\
\midrule
1 & 15/15 & 0/46 \\
2 & 15/15 & 0/46 \\
3 & 15/15 & 0/46 \\
\bottomrule
\end{tabular}
\end{table}

These are repeatability checks on the same cases, not independent
respondent samples. The zero count concerns SRE HIGH entries, not every
other cause of review. Queue creation does not establish reviewer
authorization, review completion, or correctness of human decisions.

\section*{C.1 --- Worked Score Example}

For ISCO 2211, ISIC section Q, and ISCED level 7, both implemented checks
pass. The professionals' major-group education range is $[6,8]$, with a
stricter minimum of 7 for sub-major 22. Hence
\[
  r=0.55(1)+0.45(1)=1,\qquad m=0,\qquad
  s=\min\{1,r+0.15m(1-r)\}=1.
\]
The engine marks the combination coherent and returns an adjustment of
$+0.10$. This calculation does not establish that every medical occupation
requires a master's degree.

\section*{C.2 --- Explanation Provenance}

The implementation produces English and Arabic explanations. Its
strong-alignment template currently attributes consistency to
\emph{ILO ISCO-ISIC-ISCED crosswalk tables}. That attribution does not match
the project's hand-built tables and remains an unresolved product wording
issue. The explanation should identify project-defined plausibility
rules and their limits; the thesis does not endorse an official
correspondence claim.

\chapter{Participant Information and Study Protocol}
\label{app:consent}

This material describes a proposed study. The repository submission log
records no submitted ethics application or completed recruitment.
No participant results are reported. The draft remains subject to
supervisor and institutional review.

\section*{D.1 --- Draft Participant Information}

\noindent\textbf{Study title:}
Evaluation of a Multilingual Conversational AI System for Labour Force
Survey Administration.

\noindent\textbf{Principal investigator:}
Sivarama Krishnan Subbiah (2024MCS110007), M.Tech AI and Data Science,
IIIT Kottayam.

\noindent\textbf{Supervisor:}
Dr.\ Goutam Mali, IIIT Kottayam.

The proposed study compares a conversational AI interview with an
interviewer-administered survey. Participation would involve employment
and education questions followed by a satisfaction questionnaire.
Intended languages are English, Arabic, Urdu, Hindi, and Tagalog;
translated information and consent materials are not yet finalized.

Participation would be voluntary. Approved materials must state consent,
withdrawal, retention, and deletion procedures before recruitment.
Responses would be used for research rather than government, benefit,
visa, or employment decisions. Automated codes may be inaccurate and are
not official classifications of the participant.

The data-management design must distinguish contact identifiers from
responses, define access, and explain whether text is sent to external
model providers. The prototype links identifiers to responses and can use
cloud fallback providers. Pseudonymisation is a proposed study requirement,
not an implemented application guarantee. Logging tables do not replace
consent or an approved retention plan.

\section*{D.2 --- Readiness for the Proposed Study}

\begin{table}[htbp]
\centering
\caption{Documented status and remaining requirements for participant use.}
\label{tab:app-ethics}
\small
\begin{tabularx}{\textwidth}{@{}>{\raggedright\arraybackslash}p{0.27\textwidth}p{0.22\textwidth}X@{}}
\toprule
\textbf{Requirement} & \textbf{Recorded status} & \textbf{Remaining work} \\
\midrule
Ethics application & Not submitted &
  Complete institutional forms and obtain review \\
Recruitment & Not started &
  Finalize approved eligibility and process \\
Information and consent & Draft &
  Finalize and verify translations \\
Pseudonymisation & Proposed &
  Separate identifiers and define access \\
Cloud processing & Configurable &
  Specify providers and participant disclosure \\
Supervisor authorization & Incomplete &
  Restrict review endpoints to authorized reviewers \\
Retention and withdrawal & Plan incomplete &
  Define durations, deletion, and purge operation \\
Multilingual workflow & Implemented with defects &
  Correct routing and validate supported languages \\
Pilot execution & Not started &
  Proceed after approval and readiness \\
\bottomrule
\end{tabularx}
\end{table}

The $n=30$ parallel-arm proposal is a feasibility pilot, distinct from
the English-only synthetic operational run in
Section~\ref{sec:res-synthetic-pilot}.

\chapter{Project Structure}
\label{app:structure}

This tree lists principal source areas, not every file. Agent classes,
deterministic helpers, and experimental orchestration are distinct.

\begin{lstlisting}[caption={Principal implementation and evaluation source areas.}]
.
|-- requirements.txt
|-- requirements-dev.txt
|-- alembic.ini
|-- Dockerfile
|-- docker/docker-compose.yml
|-- start.bat
|-- backend/
|   |-- main.py
|   |-- api/
|   |   |-- auth_routes.py
|   |   |-- survey_routes.py
|   |-- auth/
|   |   |-- jwt_handler.py
|   |   |-- email_otp.py
|   |-- database/
|   |   |-- models.py
|   |   |-- connection.py
|   |   |-- migrations/
|   |-- agents/
|   |   |-- conversation_manager.py
|   |   |-- language_processor.py
|   |   |-- isco_classifier.py
|   |   |-- isic_classifier.py
|   |   |-- isced_classifier.py
|   |   |-- semantic_relation.py
|   |   |-- validation_agent.py
|   |   |-- hitl_quality_manager.py
|   |   |-- report_generator.py
|   |   |-- context_memory.py
|   |   |-- audit_logger.py
|   |   |-- emotional_intelligence.py
|   |   |-- nationality_classifier.py
|   |   |-- person_register.py
|   |   |-- cross_standard_coordinator.py
|   |   |-- query_planner.py
|   |   |-- hierarchical_classification_crew.py
|   |   |-- rag_expert.py
|   |-- llm/llm_client.py
|   |-- rag/
|   |   |-- hierarchical_store.py
|   |   |-- vector_store.py
|   |   |-- official_isco08_catalogue.py
|   |   |-- standard_hierarchical_store.py
|   |   |-- standard_flat_store.py
|   |-- tests/
|-- frontend/pages/
|   |-- index.js
|   |-- chat.js
|   |-- report.js
|   |-- supervisor_review.js
|-- eval/
    |-- run_eval.py
    |-- analyze.py
    |-- build_wisco_isco_benchmark_v2_group_split.py
    |-- build_wisco_isco_benchmark_v3_dev_validation_split.py
    |-- generate_synthetic_isic_iscedf_benchmark.py
    |-- validate_synthetic_benchmark_quality.py
    |-- sre_expanded_validation.py
    |-- sre_hitl_enforcement_endpoint_verification.py
    |-- legacy_thesis_ch6/
\end{lstlisting}

The live survey invokes components directly.
\path{hierarchical_classification_crew.py} separately implements
manager-led hierarchical CrewAI delegation for experimental use and is
not imported by that route. Query planning is optional;
\texttt{RAGExpert} is a reusable helper rather than live orchestration.
Figure~\ref{fig:system-architecture} distinguishes these paths.

\chapter{Experimental Configuration and Artifact Provenance}
\label{app:reproducibility}

This appendix identifies the preserved records supporting the principal
occupation comparisons and synthetic generation. Table~\ref{tab:app-run-config}
uses short run labels for cross-reference. All five occupation runs evaluate
the same 18,747 heldout case identifiers. A profile identifies a catalogue
and encoder configuration; it is distinct from a CSV configuration hash,
an execution record, or a repository commit marker.

\section{Occupation Retrieval Configurations}

\begin{table}[htbp]
\centering
\caption{Principal occupation-run configurations and their evidence sources.
Encoder identities for B and D require profile/build evidence; the reduced
or inconsistent CSV metadata alone is insufficient.}
\label{tab:app-run-config}
\footnotesize
\setlength{\tabcolsep}{4pt}
\begin{tabularx}{\textwidth}{@{}>{\raggedright\arraybackslash}p{1.0cm}
  >{\raggedright\arraybackslash}X
  >{\raggedright\arraybackslash}p{1.9cm}
  >{\raggedright\arraybackslash}p{3.1cm}@{}}
\toprule
\textbf{Run} & \textbf{Profile and catalogue text} &
\textbf{Encoder / dimensions} & \textbf{Search and record} \\
\midrule
A-flat & \path{official_ilo2021_v1}; title-only &
E5-small / 384 & Flat; full evaluator CSV \\
A-hier & \path{official_ilo2021_v1}; title-only &
E5-small / 384 & Hierarchical; full evaluator CSV \\
B & \path{official_ilo2021_v1_e5large}; title-only &
E5-large / 1024 & Flat; reduced CSV and collection-build manifest \\
C & \path{official_ilo2021_v1_enriched}; enriched &
E5-small / 384 & Flat; full CSV and later run manifest \\
D & \path{official_ilo2021_v1_enriched_e5large}; enriched &
E5-large / 1024 & Flat; full CSV and later run manifest \\
P & \path{official_ilo2021_v1_enriched}; enriched fragments &
E5-small / 384 & Fragment-to-parent; hashed audit summary \\
\bottomrule
\end{tabularx}
\end{table}

The encoders are \path{intfloat/multilingual-e5-small} and
\path{intfloat/multilingual-e5-large}. Enriched text adds official
definitions and included occupations to the title representation.
Each flat collection contains 436 unit groups. The A-flat/A-hier command
record selects the official profile and disables LLM reranking; their CSVs
record no reranker execution. The C/D manifests also explicitly record
\texttt{use\_llm\_reranker=false}. B is described as retrieval-only in the
archived experiment summary, but its reduced CSV does not independently
record a reranker flag.

Run P is the configuration served by the local application and reported in
Section~\ref{sec:res-parent-document}. It differs from A--D in both
retrieval unit and record type. Its index holds 2,833 fragment and parent
vectors over the same 436 unit groups; a parent's score is the maximum
over its fragments, with fragment weight $0.5$. Its preserved audit record
binds the catalogue digest, fragment-map digest, selection digest, encoder
identity \path{intfloat/multilingual-e5-small}, encoder revision
\texttt{614241f6}, and an encoder weight digest, so the serving
configuration is identifiable in a way runs B and D are not. Unlike A, C,
and D, its per-case predictions are not released with this thesis; the
reported counts come from that hashed summary. Its scores are an
uncalibrated similarity blend and are not probabilities.

For A-flat/A-hier, the command record and source at the CSV-recorded
commit \texttt{280ad76} establish the following settings. The evaluator
requests five unit-group hits. A-flat returns the highest-similarity code
and exposes up to three candidates. A-hier uses beam width two:
up to two children are retained per parent at each non-leaf expansion,
rather than a globally pruned two-path frontier. Its default stage-one
mode is \texttt{description}, with a keyword major-group hint taking
precedence when available. Recorded stage-one sources are 4,100 keyword
hints and 14,647 semantic searches.

A-hier queries up to five unit-group hits per surviving minor-group
branch and selects the path whose highest-ranked leaf has the greatest
raw similarity. The weighted ancestor/leaf score supplies confidence,
not the path-selection objective. With branch collapse disabled, leaf
candidates are deduplicated across explored branches and ordered by
their highest raw similarity; up to five are exposed as candidates.
The strict hierarchical command requires genuine hierarchical execution
and sets a 30,000 ms budget per stage, shared across that stage's queries.
The saved telemetry and eligibility analysis report no fallback or
stage-budget failure in this completed run.

The historical A implementation prefixes queries with \texttt{query:},
uses normalized embeddings and query batch size one, and indexes
\texttt{passage:}-prefixed catalogue text with cosine distance.
The corresponding C source, identified by its recorded commit
\texttt{3d10925}, likewise requests five flat hits and exposes three
candidates. B's reduced export does not preserve this complete
parameter record. D's later manifest identifies flat retrieval, its
profile, and disabled reranking, but does not capture every search
parameter. Missing historical settings should therefore not be filled
from the interactive application's current defaults.

\section{Historical Records and Source Identity}

The principal artifact locations are:
\begin{itemize}[leftmargin=*]
  \item A-flat:
  \path{eval/local_runs/wisco_official_tier1_precise_deadline_full_20260809T190821Z/flat/20260809T191133Z_wisco_official_tier1_precise_deadline_full_flat.csv}.
  \item A-hier:
  \path{eval/local_runs/wisco_official_tier1_precise_deadline_full_20260809T190821Z/hierarchical/20260809T192256Z_wisco_official_tier1_precise_deadline_full_hierarchical.csv}.
  Commands are preserved in
  \path{Documentation/AI_HANDOFF/CLAUDE_TASK_36_FINAL_REPORT.md};
  the input hashes and eligibility checks are in
  \path{eval/local_runs/wisco_official_tier1_precise_deadline_full_20260809T190821Z/analysis_task37_1_clean_reproduction/analysis_manifest.json}.
  \item B:
  \path{eval/local_runs/e5large_full_heldout_20260824/results.csv},
  with encoder identity, dimensions, and verified collection counts in
  \path{eval/local_runs/e5large_build_manifest_20260823.json}.
  The build manifest establishes collection construction, not every
  setting used during the subsequent evaluation.
  \item C:
  \path{eval/results/raw_runs/enriched_catalogue_heldout_20260824/20260824T113739Z_flat.csv},
  with the later manifest
  \path{eval/results/raw_runs/enriched_catalogue_heldout_20260824/manifests/manifest_enriched_catalogue_heldout_20260824.jsonl}.
  \item D:
  \path{eval/results/raw_runs/enriched_e5large_heldout_20260824/20260824T123941Z_flat.csv},
  with the later manifest
  \path{eval/results/raw_runs/enriched_e5large_heldout_20260824/manifests/manifest_enriched_e5large_heldout_20260824.jsonl}.
\end{itemize}

A-flat/A-hier CSVs record \texttt{280ad76}; C records
\texttt{3d10925}; D records \texttt{b43f845}. The C/D manifests were
generated on 25 August 2026 and record \texttt{48f8df6}. These are separate
historical provenance fields, not the source revision of the present
thesis. B's reduced CSV stores no commit identifier.

D's CSV incorrectly names the small encoder, while its method label
and later manifest select the enriched large-encoder profile.
Moreover, the profile is absent from the catalogue registry at the
CSV-recorded commit \texttt{b43f845}. That commit marker consequently
does not establish a complete snapshot of the code executed for D;
the preserved records do not resolve the discrepancy. A subsequent
replication should archive the complete executed source, resolved encoder
revision, catalogue hashes, dependencies, and command alongside predictions.

\section{Recorded Computational Environment}

The C/D manifests describe Windows 10, Python 3.11.9, an Intel processor
reported as \emph{Intel64 Family 6 Model 154 Stepping 4, GenuineIntel},
12 logical processors, and 7.71 GiB RAM. No CUDA device is detected.
Recorded libraries include sentence-transformers 5.2.3, PyTorch 2.10.0,
transformers 4.46.3, and qdrant-client 1.17.0. The processor's commercial
model name is not recorded.

These manifests were generated after the executions and describe their
archived environment record. A and B lack an equivalent complete
environment manifest. Their hardware and dependency configuration
cannot be assumed identical to C/D merely because the experiments were
conducted on a local machine. Cross-run latency differences are descriptive
measurements, not a fully controlled hardware comparison.

\section{Synthetic-Generation Record}

The 449-row source export
\path{eval/results/synthetic_isic_iscedf_benchmark/synthetic_isic_iscedf_benchmark_20260827T021738Z.csv}
records \path{groq/openai/gpt-oss-120b} as
\texttt{generation\_model} and
\path{v2_2026-08-27_multilingual} as \texttt{prompt\_version}
for every row. It also stores the intended category and a source-definition
hash. The corresponding measured comparison is
\path{eval/results/synthetic_isic_iscedf_benchmark/synthetic_eval_results_20260827T031325Z.csv}.

The examined generator,
\path{eval/generate_synthetic_isic_iscedf_benchmark.py},
constructs the selected model through \texttt{get\_llm\_strict()} at
temperature 0.8. This is a source-configuration observation; the generation
CSV does not record temperature, sampling settings, provider model revision,
or per-call request metadata. Its model and prompt strings should not
therefore be treated as a complete historical execution trace.
The two refusal descriptions in the original measured corpus and the
distinction between generating labels and independent expert labels remain
as described in Chapters~\ref{chap:evaluation} and~\ref{chap:results}.
